\documentclass{article}
\usepackage{mathtools} 
\usepackage{fontspec}
\usepackage[UTF8]{ctex}
\usepackage{amsthm}
\usepackage{mdframed}
\usepackage{xcolor}
\usepackage{amssymb}
\usepackage{amsmath}
\usepackage{tikz}
\usepackage{pgfplots}
\usepgfplotslibrary{polar}
\usetikzlibrary{3d} % 允许3D坐标
% \pgfplotsset{compat=1.18}


% 定义新的带灰色背景的说明环境 zremark
\newmdtheoremenv[
  backgroundcolor=gray!10,
  % 边框与背景一致，边框线会消失
  linecolor=gray!10
]{zremark}{说明}

% 通用矩阵命令: \flexmatrix{矩阵名}{元素符号}{行数}{列数}
\newcommand{\flexmatrix}[4]{
  \[
  #1 = \begin{pmatrix}
    #2_{11}     & #2_{12}     & \cdots & #2_{1#4}   \\
    #2_{21}     & #2_{22}     & \cdots & #2_{2#4}   \\
    \vdots      & \vdots      & \ddots & \vdots     \\
    #2_{#31}    & #2_{#32}    & \cdots & #2_{#3#4}
  \end{pmatrix}
  \]
}

% 简化版命令（默认矩阵名为A，元素符号为a）: \quickmatrix{行数}{列数}
\newcommand{\quickmatrix}[2]{\flexmatrix{A}{a}{#1}{#2}}

\begin{document}
\title{习题 11.1}
\author{张志聪}
\maketitle

\section*{1}

\begin{itemize}
  \item (1)
        \begin{align*}
          \int_{0}^{+\infty} x e^{-x^2} dx
           & = \lim\limits_{u \to +\infty} \int_{0}^{u} x e^{-x^2} dx            \\
           & = \lim\limits_{u \to +\infty} - \frac{1}{2}\int_{0}^{u} d(e^{-x^2}) \\
           & = \lim\limits_{u \to +\infty} - \frac{1}{2} (e^{-u^2} - 1)          \\
           & = \frac{1}{2}
        \end{align*}
        所以，收敛。
  \item (8)

        利用第八章$\S 2$例8
        \begin{align*}
          \int_{0}^{+\infty} \frac{dx}{\sqrt{1 + x^2}}
           & = \lim\limits_{u \to +\infty} \int_{0}^{u} \frac{dx}{\sqrt{1 + x^2}}   \\
           & = \lim\limits_{u \to +\infty} (\ln |x + \sqrt{1 + x^2}|)\bigg|_{0}^{u} \\
           & = \lim\limits_{u \to +\infty} \ln |u + \sqrt{1 + u^2}|                 \\
           & = + \infty
        \end{align*}
        所以，发散。
\end{itemize}

\section*{2}

\begin{itemize}
  \item (1)

        令$t = x - a, dx = dt, x: a \to b, t: 0 \to b - a$，
        我们有
        \begin{align*}
          \int_{a}^{b} \frac{dx}{(x - a)^p}
           & = \int_{0}^{b - a} \frac{dt}{t^p}
        \end{align*}
        点$0$是被积函数$\frac{1}{t^p}$的瑕点，且在$(0, b-a]$上连续，从而在任何$[0, u] \subset (0, b-a]$上可积，
        依定义
        \begin{align*}
          \int_{0}^{b - a} \frac{dt}{t^p}
           & = \lim\limits_{u \to 0^{+}} \frac{dt}{t^p}
        \end{align*}
        由例6可知，$0 < q < 1$时，瑕积分收敛，且
        \begin{align*}
          \lim\limits_{u \to 0^{+}} \frac{dt}{t^p} = \frac{(b - a)^{1 - q}}{1-q}
        \end{align*}
        而当$q \geq 1$时，瑕积分收敛于$+\infty$。

        另外，$p \leq 0$时，$-p \geq 0$，则$\frac{1}{(x - a)^p} = (x - a)^{-p}$没有瑕点，
        题目变为一般定积分，不做讨论。

  \item (8)

        \begin{itemize}
          \item $p = 1$；

                $x = 0, 1$都是函数$f(x) = \frac{1}{x(\ln x)^p}$的瑕点。
                由定义可知
                \begin{align*}
                  \int_{0}^{1} \frac{dx}{x\ln x}
                   & = \int_{0}^{\frac{1}{2}} \frac{dx}{x\ln x} + \int_{\frac{1}{2}}^{1} \frac{dx}{x\ln x}                                                               \\
                   & = \lim\limits_{u \to 0^{+}} \int_{u}^{\frac{1}{2}} \frac{dx}{x\ln x} + \lim\limits_{v \to 1^{-}} \int_{\frac{1}{2}}^{v} \frac{dx}{x\ln x}           \\
                   & = \lim\limits_{u \to 0^{+}} \int_{u}^{\frac{1}{2}} \frac{d(\ln x)}{\ln x} + \lim\limits_{v \to 1^{-}} \int_{\frac{1}{2}}^{v} \frac{d(\ln x)}{\ln x} \\
                   & = \lim\limits_{u \to 0^{+}} \ln |\ln x| \bigg|_u^{\frac{1}{2}} + \lim\limits_{v \to 1^{-}} \ln |\ln x| \bigg|_v^{\frac{1}{2}}                       \\
                   & = \lim\limits_{u \to 0^{+}} \ln |\ln \frac{1}{2}| - \ln |\ln u| + \lim\limits_{v \to 1^{-}} \ln |\ln v| - \ln |\ln \frac{1}{2}|
                \end{align*}
                第一个极限$-\infty$，所以发散。

          \item $p \neq 1$；

                $x = 0, 1$都是函数$f(x) = \frac{1}{x(\ln x)^p}$的瑕点。
                由定义可知
                \begin{align*}
                  \int_{0}^{1} \frac{dx}{x(\ln x)^p}
                   & = \int_{0}^{\frac{1}{2}} \frac{dx}{x(\ln x)^p} + \int_{\frac{1}{2}}^{1} \frac{dx}{x(\ln x)^p}       \\
                   & = \lim\limits_{u \to 0^{+}} \frac{dx}{x(\ln x)^p} + \lim\limits_{v \to 1^{-}} \frac{dx}{x(\ln x)^p} \\
                   & = \lim\limits_{u \to 0^{+}} \frac{1}{1 - p}(\ln x)^{1-p}\bigg|_{u}^{\frac{1}{2}}
                  + \lim\limits_{v \to 1^{-}} \frac{1}{1 - p}(\ln x)^{1-p}\bigg|_{\frac{1}{2}}^{v}                       \\
                   & = \lim\limits_{u \to 0^{+}} \frac{1}{1 - p}[(\ln \frac{1}{2})^{1-p} - (\ln u)^{1-p}]
                  + \lim\limits_{v \to 1^{-}} \frac{1}{1 - p}[(\ln v)^{1-p} - (\ln \frac{1}{2})^{1-p}]
                \end{align*}

                当$0 < p < 1$时，$1 - p > 0$，第一个极限为发散。

                当$1 < p$时，第二个极限发散，故发散。
        \end{itemize}

        综上，发散。

        \textbf{注：} 题目本身已经确保$(ln x)^p$是有意义的，这使得$(\ln x)^{1-p}$也是有意义的，
        否则$\ln x$可能是负数，而$1 - p = \frac{1}{2}$时，出现矛盾。

\end{itemize}

\section*{3}

利用例6，令
\begin{align*}
  f(x)   & = \frac{1}{x^{\frac{1}{2}}} = \frac{1}{\sqrt{x}} \\
  f^2(x) & = \frac{1}{x}
\end{align*}
于是
\begin{align*}
  \int_{0}^{1} f(x)
\end{align*}
收敛
\begin{align*}
  \int_{0}^{1} f^2(x)
\end{align*}
发散。

\section*{4}

出现这种现象的原因：

设$\int_{a}^{+\infty} f(x) dx$收敛，于是我们有
\begin{align*}
  \int_{a}^{+\infty} f(x) dx & = \lim\limits_{u \to +\infty} \int_{a}^{u} f(x) dx                        \\
                             & = \lim\limits_{u \to +\infty} \int_{a}^{k} f(x) dx + \int_{k}^{u} f(x) dx \\
                             & = \int_{a}^{k} f(x) dx + \lim\limits_{u \to +\infty} \int_{k}^{u} f(x) dx
\end{align*}
考虑上式右侧，令$k \to +\infty$，按照无穷积分的定义
\begin{align*}
  \lim \limits_{k \to +\infty} \int_{a}^{k} f(x) dx = \int_{a}^{+\infty} f(x) dx
\end{align*}
于是
\begin{align*}
  \lim\limits_{k \to +\infty} \lim\limits_{u \to +\infty} \int_{k}^{u} f(x) dx = 0
\end{align*}
按照极限定义，对任意$\epsilon > 0$，存在$M > 0$，当$k > M$时，我们有
\begin{align*}
  \left|\lim\limits_{u \to +\infty} \int_{k}^{u} f(x) dx \right| < \epsilon
\end{align*}
这只能得到$f(x)$这区间$(k, +\infty)$上的面积$< \epsilon$，但与$f(x)$的值是否
小于$\epsilon$无关。

\section*{5}

若$\lim\limits_{x \to +\infty} f(x) dx = A$，且$A > 0$，则存在$M > 0$，
有$x \geq M$，有$f(x) > \frac{A}{2}$。于是对任意$u > M$，有
\begin{align*}
  \int_{a}^{u} f(x) dx & = \int_{a}^{M} f(x) dx + \int_{M}^{u} f(x) dx                                        \\
                       & \geq  \int_{a}^{M} f(x) dx + (u - M)\cdot \frac{A}{2} \to +\infty \, (u \to +\infty)
\end{align*}
这与$\int_{a}^{+\infty} f(x) dx$收敛相矛盾。

同理可证$A < 0$。

综上，$A = 0$。

\section*{6}

$\int_{a}^{+\infty} f'(x) dx$收敛，即
\begin{align*}
  \int_{a}^{+\infty} f'(x) dx 
  & = \lim\limits_{u \to +\infty} \int_{a}^{u} f'(x) dx \\
  & = \lim\limits_{u \to +\infty} f(u) - f(a) 
\end{align*}
收敛，于是可得$f(x)$收敛，由习题5可得$\lim\limits_{x \to +\infty} f(x) = 0$。


\end{document}